\section{Related Works}
\label{sec:rel}

\paragraph{Distributed Storage Systems.}
As mentioned in \cref{sec:motivation}, our work is related to a long line of research and production systems in distributed storage~\cite{ceph,gfs,bigtable,dynamo,liquid,amazon-s3,amazon-glacier,google-storage,google-cloudstore,google-coldline}.
Unlike \sys, all those systems are centrally managed by a single administrative entity, and all participants are assumed to be non-Byzantine.
Our object store interface is similar to Amazon S3~\cite{amazon-s3} and Google Cloud Storage~\cite{google-cloudstore}.
Many prior systems deploy a centralized service to manage storage metadata~\cite{gfs,bigtable}, while metadata management is fully decentralized in \sys.
Dynamo~\cite{dynamo} uses consistent hashing~\cite{chash} to assign keys to nodes, similar to our DHT-based approach.
Object-to-server mapping in Ceph~\cite{ceph} is done through a distribution function CRUSH~\cite{crush} which maps each object to a placement group based on the object hash.
\ifspace
Our approach of using rateless erasure code is inspired by liquid storage~\cite{liquid}.
Other erasure codes such as local reconstruction codes~\cite{lrcode}, regenerating codes~\cite{rcode}, and traditional MDS codes have also been used in distributed storage.
\fi

\paragraph{Decentralized Storage.}
Building reliable storage systems in a fully decentralized environment has been explored extensively~\cite{xfs,farsite,oceanstore,filecoin,arweave}.
Farsite~\cite{farsite} uses BFT replication for directory metadata and CFT replication for file data.
Membership management is done through trusted certificate authorities.
\sys reduces storage redundancy by using erasure coding, and avoids centralized trust through PoS and verifiable random peer selection.
Most closely related to our work is the deep archival storage in OceanStore which stores erasure coded fragments over multiple failure domains.
However, OceanStore makes centralized placement decisions and relies on trusted, centralized nodes to perform repairs.
As detailed in \cref{sec:motivation}, IPFS~\cite{filecoin} uses DHT to store publisher record without centralization, but offers no mechanism to reliably store objects.
Filecoin~\cite{filecoin} and Arweave~\cite{arweave} both use tokens to incentivize nodes to store objects longer, but has weak durability guarantees in the presence of failures.

\ifspace
\paragraph{BFT Protocols.}
There exists a long line of research on replication protocols that tolerant Byzantine failures~\cite{pbft,algorand,hotstuff,bitcoin}.
These protocols, however, require data to be replicated on all participants, sacrificing storage efficiency.
Our peer selection protocol using verifiable random function is inspired by Algorand~\cite{algorand} and Sleepy Consensus~\cite{sleepyconsensus}.
Our proof-of-stake based approach to defend Sybil attacks is similar to the ones used in Ethereum~\cite{ethereum} and Algorand~\cite{algorand}.
\fi

\iffalse

\textbf{Rateless erasure coding.} RaptorQ (IETF spec) performs well with large range of $k$.

\textbf{Gossip protocols.} Floodsub by libp2p.
Gossipsub by libp2p limits maximum outbound degree to avoid congestion on dense-connected peers.
Kadcast selects mesh from Kademlia buckets.

\textbf{Other usage of VRF.} Algorand.

\textbf{Distributed algorithm for locating data objects.} CRUSH by Ceph.

\sgd{Include this if we have discussion part.} \textbf{Aggressive attacking.}
In this work we mainly discuss on how to achieve long-term unpredictable;
however, in short term attacker may predict the next responsible peers by observing peers pulling fragment (which involves exchanging Kleroterion proof which attacker can also verify).
If the attacker can compromise the peer before it become responsible and send out first gossip, the safety invariant breaks.
Notice that the lazy fetching optimization previously described in \XXX can effectively shorten attacking time window.
With active redundancy reduction (\autoref{sec:implementation}), the attacking is also mitigated.
There could be more than enough number of nodes fetching fragments, where many of them silently discard the fragment later and make the attacking meaningless.
This will increase the cost of a successful attacking significantly.

\fi